% Generated draft: requires mathematical and source review.
\documentclass[11pt]{article}
\usepackage[margin=25mm]{geometry}
\usepackage{fontspec}
\setmainfont{Latin Modern Roman}
\usepackage{amsmath,amssymb,mathtools,unicode-math}
\setmathfont{Latin Modern Math}
\usepackage{enumitem,needspace,xurl,hyperref}
\setlength{\parindent}{0pt}
\setlength{\parskip}{5pt}
\setlength{\emergencystretch}{3em}
\newcommand{\N}{\mathbb{N}}
\newcommand{\Z}{\mathbb{Z}}
\newcommand{\Q}{\mathbb{Q}}
\newcommand{\R}{\mathbb{R}}
\newcommand{\C}{\mathbb{C}}
\newcommand{\F}{\mathbb{F}}
\newcommand{\A}{\mathbb{A}}
\newcommand{\PP}{\mathbb{P}}
\newcommand{\E}{\mathbb{E}}
\newcommand{\eps}{\varepsilon}
\newcommand{\dd}{\,\mathrm{d}}
\begin{document}
\section*{UnsolvedMath ID: 30002934}
\textbf{Source problem code:} catalogue entry
\subsection{Definitions and mathematical statement}
A closed hypersurface means a compact hypersurface without boundary. A minimal hypersurface has identically vanishing mean curvature vector. The Morse index \(\operatorname{index}(M)\) is understood here as the maximal dimension of a space of smooth normal variations on which the second variation of volume is negative definite. This definition uses sections of the normal bundle and does not require choosing a global unit normal; the source does not explicitly specify its index convention.

Fix a Riemannian manifold \(N^{n+1}\) with strictly positive sectional curvature, meaning \(\sec_N(\sigma)>0\) for every tangent two-plane \(\sigma\). Write \(\mathcal H^n(M)\) for the source's \(n\)-dimensional hypersurface volume. For an embedded hypersurface this is its Hausdorff measure in the ambient metric. The source does not specify how volume or multiplicity should be interpreted if immersed hypersurfaces are included.

\textbf{Conjectural statement as supplied.} For every nonnegative integer \(I\), does there exist a finite constant \(C(N,I)\) such that every closed minimal hypersurface \(M^n\) in the fixed ambient manifold \(N\) satisfies
\[
\operatorname{index}(M)=I
\quad\Longrightarrow\quad
\mathcal H^n(M)\leq C(N,I)?
\]
The constant must be independent of \(M\). The source also asks for the resulting sequence formulations: for every sequence \((M_j)_{j\geq1}\) of such hypersurfaces in the same \(N\),
\[
\sup_{j\geq1}\operatorname{index}(M_j)<\infty
\quad\Longrightarrow\quad
\sup_{j\geq1}\mathcal H^n(M_j)<\infty,
\]
and
\[
\mathcal H^n(M_j)\longrightarrow\infty
\quad\Longrightarrow\quad
\operatorname{index}(M_j)\longrightarrow\infty.
\]

\textbf{Unresolved scope.} The supplied formulation does not specify a dimension range, whether \(N\) must be compact or complete, whether hypersurfaces must be embedded or may be immersed, whether they must be connected, or whether orientability or two-sidedness is required. Strict positivity of sectional curvature is stated pointwise; no uniform positive lower bound is supplied. These omissions leave the intended class of hypersurfaces unresolved, so the displayed statement requires review against the original passage before being treated as a fully specified conjecture.
\subsection{Short English statement}
The supplied conjecture asks whether the volume of a closed minimal hypersurface in a fixed positively curved ambient manifold is bounded solely by its Morse index and the ambient manifold. It also asks whether bounded index forces uniformly bounded volume, and whether volume tending to infinity forces index tending to infinity. The intended geometric class remains incompletely specified.
\subsection{Sources}
\begin{itemize}
\item Local UnsolvedMath registry, supplied complete source record: permanent ID \(30002934\), exact source problem code \texttt{catalogue entry}, title \emph{Index Bounds for Minimal-Hypersurface Volume}, in the Oberwolfach Reports Open Problems collection. The record gives proposed year \(2015\) and supplies no named proposer.
\item Supplied report citation: \emph{Partial Differential Equations (2016)}, Oberwolfach Reports, pp.~2065--2124, DOI \texttt{10.4171/owr/2015/36}. \url{https://doi.org/10.4171/owr/2015/36}. The background describes Oberwolfach Reports as documenting mathematical workshops and their open-problem discussions. The supplied title year \(2016\) and the \(2015\) year appearing elsewhere in the record are retained without reconciliation.
\item Supplied literature source: Buzano--Sharp, ``bounded index and area estimates.'' \url{https://arxiv.org/abs/1512.01047}.
\item Supplied literature source: ``Oberwolfach Report on minimal hypersurfaces.'' \url{https://doi.org/10.4171/owr/2015/36}.
\item Literature assessment dated exactly \texttt{2026-08-21}: the supplied curation labels the problem open and states that bounded-index compactness and curvature quantization are known with an area bound, that the surface case has stronger positive-curvature results, and that the proposed uniform area bound from index and positive sectional curvature remains conjectural in general dimensions. This is supplied dated curation derived from the OWR record, not an independently verified literature assessment; its status and cited claims require independent verification.
\item The registry labels the statement \texttt{corrected\_verified} and reports that the official Oberwolfach passage was checked directly and its question, definitions, and notation restored without strengthening the claim. This is a supplied verification claim, not verification performed for this draft.
\end{itemize}
\end{document}
